%%%%%%%%%%%%%%%%%%%%%%%%%%%%%%%%%%%%%%%%%%%%%%%%%%%%%%%%%%%%%%%%%%%%%%%%%%%%%%%
% CHAPTER: ELLIPTIC ALGEBRAS AND THE K3 CHIRAL ALGEBRA
%%%%%%%%%%%%%%%%%%%%%%%%%%%%%%%%%%%%%%%%%%%%%%%%%%%%%%%%%%%%%%%%%%%%%%%%%%%%%%%

\chapter{Elliptic algebras and the K3 chiral algebra}
\label{chap:toroidal-elliptic}

The elliptic curve $E_\tau$ plays two roles in this chapter: as the base curve
carrying chiral algebras (toroidal, elliptic quantum groups, Felder's dynamical
$R$-matrix) and as a geometric factor in the Calabi--Yau threefold $K3 \times E$
whose BPS spectrum is governed by the Igusa cusp form $\Phi_{10}$.
The link between these roles is the Hilbert scheme $\operatorname{Hilb}^n(K3)$:
the toroidal algebra $U_{q,t}(\hat{\hat{\mathfrak{g}}})$ acts on
its $K$-theory (Schiffmann--Vasserot), while the DMVV formula
$1/\Phi_{10} = \sum_N p^N\,\chi(\operatorname{Sym}^N(K3))$
encodes the character.
The Fourier--Jacobi coefficients of $\Phi_{10}$ are Jacobi forms on $E_\tau$
built from the same theta functions that define the elliptic propagator.


%% ===============================================================
%% ACT I: THE ELLIPTIC PROPAGATOR
%% ===============================================================

\section{The elliptic propagator}
\label{sec:elliptic-propagator}

\subsection{From Arnold to Fay}
\label{subsec:arnold-to-fay}

\begin{definition}[Elliptic propagator]
\label{def:elliptic-propagator}
\index{elliptic propagator|textbf}
\index{Weierstrass zeta function}
On the elliptic curve $E_\tau = \bC/(\bZ + \tau\bZ)$,
the \emph{elliptic propagator} is
\begin{equation}\label{eq:ell-propagator}
\eta_{ij}^{\mathrm{ell}}
= d\log\theta_1(z_i - z_j \,|\, \tau)
= \zeta_\tau(z_i - z_j)\,(dz_i - dz_j),
\end{equation}
where $\zeta_\tau(z) = \partial_z \log\theta_1(z|\tau)$
is the Weierstrass zeta function on $E_\tau$.  This
replaces the rational propagator $\eta_{ij}^{\mathrm{rat}}
= d\log(z_i - z_j) = dz_{ij}/(z_i - z_j)$ on $\bC$.
\end{definition}

The Laurent expansion of $\zeta_\tau$ about $z = 0$
encodes the Eisenstein series:
\begin{equation}\label{eq:zeta-laurent}
\zeta_\tau(z)
= \frac{1}{z}
- \frac{\pi^2}{3}\,E_2(\tau)\,z
- \frac{\pi^4}{45}\,E_4(\tau)\,z^3
- \frac{2\pi^6}{945}\,E_6(\tau)\,z^5
- \cdots
\end{equation}
The leading term $1/z$ is the rational propagator; the
correction $\delta(\tau) := \zeta_\tau(z) - 1/z$ is regular
and encodes the genus-$1$ geometry of $E_\tau$ through
(quasi-)modular forms.

\begin{definition}[Jacobi theta function]\label{def:theta}
\index{Jacobi theta function}
The Jacobi theta function is
\[
\theta(u \,|\, \tau)
= \sum_{n \in \bZ} (-1)^n\, e^{\pi i \tau\, n(n-1)
+ 2\pi i n u}
= (e^{2\pi iu}; p)_\infty\,
(p\, e^{-2\pi iu}; p)_\infty\,
(p; p)_\infty
\]
where $p = e^{2\pi i \tau}$ and $(a; q)_\infty
= \prod_{n \geq 0}(1 - a\,q^n)$.
\end{definition}

\begin{theorem}[Fay trisecant identity at genus~$1$
{\cite{Fay73}}; \ClaimStatusProvedElsewhere]
\label{thm:fay}
\index{Fay trisecant identity|textbf}
For any four points $u_1, u_2, u_3, u_4$ on $E_\tau$:
\begin{equation}\label{eq:fay}
\theta(u_1{-}u_2)\,\theta(u_3{-}u_4)
- \theta(u_1{-}u_3)\,\theta(u_2{-}u_4)
+ \theta(u_1{-}u_4)\,\theta(u_2{-}u_3)
= 0.
\end{equation}
In the rational limit $\tau \to i\infty$,
$\theta_1(z|\tau) \to 2\sin(\pi z) \to 2\pi z$
and~\eqref{eq:fay} reduces to the Ptolemy relation.
\end{theorem}

The Fay identity is the genus-$1$ Arnold relation.  The
rational Arnold relation
$\eta_{12} \wedge \eta_{23} + \eta_{23} \wedge \eta_{31}
+ \eta_{31} \wedge \eta_{12} = 0$
(Theorem~\ref{thm:arnold-relations}) governs the genus-$0$
bar differential; the Fay identity governs the elliptic bar
differential.


\subsection{Bar nilpotency from Fay}
\label{subsec:bar-nilpotency-fay}

\begin{proposition}[Fay identity implies elliptic
\texorpdfstring{$d^2 = 0$}{d2 = 0};
\ClaimStatusProvedHere]
\label{prop:fay-implies-d-squared}
\index{Fay trisecant identity!and bar nilpotency}
On $\overline{C}_3(E_\tau)$, the elliptic bar differential
satisfies $d^2 = 0$.  The algebraic input is the Fay
trisecant identity (Theorem~\textup{\ref{thm:fay})},
replacing the Arnold relation in the rational case
(Theorem~\textup{\ref{thm:bar-nilpotency-complete}}).
\end{proposition}

\begin{proof}
The computation of $d^2$ on a degree-$1$ element
$[a] \otimes 1$ produces terms on
$\overline{C}_3(E_\tau)$ involving the $2$-form
\begin{equation}\label{eq:elliptic-arnold}
\eta_{12}^{\mathrm{ell}} \wedge \eta_{23}^{\mathrm{ell}}
+ \eta_{23}^{\mathrm{ell}} \wedge \eta_{31}^{\mathrm{ell}}
+ \eta_{31}^{\mathrm{ell}} \wedge \eta_{12}^{\mathrm{ell}}
= 0,
\end{equation}
the elliptic Arnold relation.  This follows from the Fay
identity by differentiating~\eqref{eq:fay} with respect
to $u_1$ and $u_4$, then setting $u_4 = u_1$.  The
resulting identity is
\begin{multline}\label{eq:fay-differentiated}
\zeta_\tau(u_1{-}u_2)\,\zeta_\tau(u_2{-}u_3)
+ \zeta_\tau(u_2{-}u_3)\,\zeta_\tau(u_3{-}u_1)
+ \zeta_\tau(u_3{-}u_1)\,\zeta_\tau(u_1{-}u_2) \\
= \wp_\tau(u_1{-}u_2) + \wp_\tau(u_2{-}u_3)
+ \wp_\tau(u_3{-}u_1),
\end{multline}
where $\wp_\tau = -\zeta_\tau'$.  Taking exterior
derivatives, the right side contributes no mixed wedge
products (each $\wp_\tau(z_i - z_j)$ depends on a single
difference), so~\eqref{eq:elliptic-arnold} follows.

The remaining terms in $d^2$ cancel by OPE associativity:
$d = d_{\mathrm{local}} + d_{\mathrm{global}}$, where
$d_{\mathrm{local}}^2 = 0$ by the OPE of $\cA$, and the
cross-terms vanish because the elliptic propagator has the
same residue structure as the rational one
($\Res_{z=0} \zeta_\tau(z)\,dz = 1$).
\end{proof}

\begin{remark}[Rational limit]\label{rem:fay-to-arnold}
As $\tau \to i\infty$: $\zeta_\tau(z) \to 1/z$,
$\wp_\tau(z) \to 1/z^2$,
and~\eqref{eq:elliptic-arnold} reduces to the Arnold
relation (Theorem~\ref{thm:arnold-relations}).
\end{remark}


\subsection{The elliptic bar complex}
\label{subsec:elliptic-bar-complex}

\begin{construction}[Elliptic bar complex]
\label{constr:elliptic-bar}
\index{bar complex!elliptic|textbf}
Replace the rational propagator $d\log(z_1 - z_2)$
with $\eta_{12}^{\mathrm{ell}}
= d\log\theta_1(z_1 - z_2 \,|\, \tau)$ in the bar
complex.  The bar differential extracts residues at
the zeros of $\theta$, which occur at the lattice
points $u \in \bZ + \tau\bZ$.
\end{construction}

\begin{theorem}[Elliptic versus rational bar homology;
\ClaimStatusProvedHere]
\label{thm:elliptic-vs-rational}
\index{bar complex!elliptic vs rational}
Let $\cA$ be a chiral algebra on $E_\tau$.  The elliptic
bar homology decomposes as
\[
H_n(\B^{\mathrm{ell}}(\cA))
\;\cong\;
H_n(\B^{\mathrm{rat}}(\cA))
\;\oplus\;
\bigoplus_{k \geq 1}
\mathcal{E}^{(k)}_n(\cA, \tau),
\]
where the Eisenstein corrections
$\mathcal{E}^{(k)}_n(\cA, \tau)$ lie in the space of
\textup{(}quasi-\textup{)}modular forms of weight~$2k$:
\begin{enumerate}[label=\textup{(\roman*)}]
\item $\mathcal{E}^{(1)}_n \propto E_2(\tau)$
  \textup{(}quasi-modular, weight~$2$\textup{)};
\item $\mathcal{E}^{(k)}_n \in
  M_{2k}(\mathrm{SL}_2(\bZ))$ for $k \geq 2$
  \textup{(}holomorphic modular\textup{)};
\item $\mathcal{E}^{(k)}_n = 0$ for $k \geq n$;
\item at $\tau \to i\infty$: all corrections degenerate
  to constants.
\end{enumerate}
\end{theorem}

\begin{proof}
Filter the elliptic bar complex by correction order,
expanding each propagator via~\eqref{eq:zeta-laurent}
as $\zeta_\tau(z) = z^{-1}
+ \sum_{k \geq 1} c_k(\tau)\,z^{2k-1}$,
where $c_1(\tau) = -(\pi^2/3)\,E_2(\tau)$ and
$c_k(\tau) \in M_{2k+2}(\mathrm{SL}_2(\bZ))$
for $k \geq 2$.  The associated spectral sequence has
$E_1^{0,n} = H_n(\B^{\mathrm{rat}}(\cA))$ and converges
(the filtration is bounded: at most $n-1$ propagators
at bar degree~$n$).  The modular weight grading
splits the filtration into a direct sum, by linear
independence of (quasi-)modular forms of distinct weights.
Zhu's modular invariance theorem~\cite{Zhu96} ensures
that genus-$1$ correlation functions transform as
vector-valued modular forms, so the differentials
respect the weight grading.
\end{proof}


\subsection{Explicit: Heisenberg bar complex on
\texorpdfstring{$E_\tau$}{E tau}}
\label{subsec:ell-bar-heisenberg}
\label{sec:elliptic-bar-heisenberg}

\begin{computation}[Heisenberg elliptic bar differential]
\label{comp:ell-bar-deg2}
\index{bar complex!elliptic!Heisenberg}
For the Heisenberg algebra $\cH_\kappa$ on $E_\tau$,
the bar differential at degree~$2$ is:
\begin{equation}\label{eq:ell-bar-diff-result}
d^{\mathrm{ell}}([a_m | a_n] \otimes
\eta_{12}^{\mathrm{ell}})
= -\frac{\kappa\pi^2}{3}\,E_2(\tau)
\cdot z_2^{-m-n-2}\,dz_2.
\end{equation}
This arises from the residue
$\Res_{\epsilon = 0}
[\kappa\,\epsilon^{-2}\,\zeta_\tau(\epsilon)\,
d\epsilon]
= -(\kappa\pi^2/3)\,E_2(\tau)$:
the $\epsilon^{-3}$ term has vanishing residue, and
the first nonzero contribution is the $E_2$ correction.
The rational bar differential ($\zeta_\tau \to 1/z$)
vanishes identically: the \emph{entire degree-$2$
bar differential is the elliptic correction}.
The curvature element is
$m_0^{\mathrm{ell}} = \kappa \cdot (\pi^2/3)\,
E_2(\tau)$, encoding the failure of double periodicity.
\end{computation}

\begin{table}[ht]
\centering
\caption{Heisenberg elliptic versus rational bar
differentials at low degrees}
\label{tab:ell-vs-rat-bar}
\begin{tabular}{|c|c|c|c|}
\hline
\textbf{Degree} & \textbf{Rational $d^{\mathrm{rat}}$}
& \textbf{Elliptic $d^{\mathrm{ell}}$}
& \textbf{Correction} \\ \hline
$1$ & $0$ & $0$ & --- \\ \hline
$2$ & $0$ & $-\frac{\kappa\pi^2}{3}\,E_2(\tau)$
& $E_2$ (wt~$2$) \\ \hline
$3$ & $0$ & $-\frac{\kappa\pi^4}{45}\,E_4(\tau)$
& $E_4$ (wt~$4$) \\ \hline
$4$ & $0$ & $-\frac{2\kappa\pi^6}{945}\,E_6(\tau)$
& $E_6$ (wt~$6$) \\ \hline
\textbf{General $n$} & $0$
& $-\kappa \cdot \frac{2^{2n-2}|B_{2n-2}|}{(2n-2)!}
  \cdot \pi^{2n-2}\,E_{2n-2}(\tau)$
& $E_{2n-2}$ (wt~$2n{-}2$) \\ \hline
\end{tabular}
\end{table}


%% ===============================================================
%% ACT II: ALGEBRAS ON THE ELLIPTIC CURVE
%% ===============================================================

\section{Algebras on the elliptic curve}
\label{sec:algebras-on-elliptic}

\subsection{The toroidal algebra}
\label{subsec:toroidal-algebra}
\label{sec:double-affine}

\begin{definition}[Toroidal algebra]
\label{def:toroidal}
\index{toroidal algebra|textbf}
\index{double affine algebra|textbf}
The \emph{toroidal algebra}
$U_{q,t}(\hat{\hat{\mathfrak{g}}})$
is generated by horizontal currents $x^\pm_{i,n}$,
$\psi^\pm_{i,n}$ (affine in one direction) and
vertical currents $e_i(z)$, $f_i(z)$, $\psi_i(z)$
(affine in the other), subject to Drinfeld-type
relations in each direction,
cross-relations
$\psi_i(z)\,e_j(w)
= g_{ij}(z/w)\,e_j(w)\,\psi_i(z)$
with $g_{ij}(u)
= (u - q^{a_{ij}})/(u - q^{-a_{ij}})$,
and Serre relations in both directions.  The two
independent parameters are $q$ and $t = q^k$.
\end{definition}

\begin{proposition}[Toroidal OPE {\cite{FBZ04}};
\ClaimStatusProvedElsewhere]
\label{prop:toroidal-ope}
\index{toroidal algebra!OPE}
The OPE of toroidal currents involves both parameters:
\[
e_i(z)\,f_j(w)
\sim \frac{\delta_{ij}}{(1 - q)(1 - t^{-1})}
\cdot \frac{\psi_i^+(w) - \psi_i^-(w)}{z - w}
+ \text{rational in } q, t.
\]
\end{proposition}

\begin{conjecture}[Elliptic-curve toroidal realization;
\ClaimStatusConjectured]
\label{conj:toroidal-e1}
\index{toroidal algebra!E1-chiral realization}
$U_{q,t}(\hat{\hat{\mathfrak{g}}})$ admits a curve-based
$\Eone$-chiral realization on $E_\tau$ with
$\tau = \log t / \log q$.
\end{conjecture}

\begin{conjecture}[Toroidal surface-factorization;
\ClaimStatusConjectured]
\label{conj:toroidal-surface-factorization}
The toroidal data extends to a doubly-graded
surface-factorization object on
$\bC^* \times \bC^*$, beyond the curve-chiral
hierarchy of this monograph.
\end{conjecture}

\begin{conjecture}[Toroidal Koszul dual;
\ClaimStatusConjectured]
\label{conj:toroidal-koszul-dual}
\label{rem:toroidal-three-theorems}%
\index{toroidal algebra!Koszul dual}
The $\Eone$-chiral Koszul dual of
$U_{q,t}(\hat{\hat{\mathfrak{g}}})$ is
$U_{q^{-1},t^{-1}}(\hat{\hat{\mathfrak{g}}})$.
The evidence: RTT inversion
$R(u;q,t)^{-1} = R(u;q^{-1},t^{-1})$,
affine degeneration at $t \to 0$, and the
$q \to q^{-1}$ reversal of complex structure on $E_\tau$.
\end{conjecture}


\subsection{Elliptic quantum groups}
\label{subsec:elliptic-quantum-groups}
\label{sec:elliptic-quantum}

\begin{definition}[Elliptic quantum group]
\label{def:elliptic-quantum}
\index{elliptic quantum group|textbf}
The \emph{elliptic quantum group}
$E_{q,p}(\mathfrak{g})$ (Felder) is defined by a
dynamical $R$-matrix $R(u, \lambda)$ depending on
spectral parameter $u$ and dynamical parameter
$\lambda \in \mathfrak{h}^*$, with the
\emph{dynamical RLL relations}:
\[
R(u{-}v, \lambda)\,L_1(u, \lambda)\,
L_2(v, \lambda - h^{(1)})
= L_2(v, \lambda)\,L_1(u, \lambda - h^{(2)})\,
R(u{-}v, \lambda).
\]
\end{definition}

\begin{definition}[Elliptic $R$-matrix for
$\mathfrak{sl}_2$]\label{def:elliptic-r}
\index{R-matrix!elliptic}
\[
R(u, \lambda)
= \begin{pmatrix}
a(u) & 0 & 0 & 0 \\
0 & b(u, \lambda) & c(u, \lambda) & 0 \\
0 & c(u, -\lambda) & b(u, -\lambda) & 0 \\
0 & 0 & 0 & a(u)
\end{pmatrix}
\]
with
$a(u) = \theta(u{+}\eta)/\theta(\eta)$,\quad
$b(u, \lambda)
= \theta(u)\,\theta(\lambda{+}\eta)
/(\theta(\eta)\,\theta(\lambda))$,\quad
$c(u, \lambda)
= \theta(u{+}\lambda)\,\theta(\eta)
/(\theta(\lambda)\,\theta(u{+}\eta))$.
\end{definition}


\subsection{Dynamical Yang--Baxter equation}
\label{subsec:dybe}
\label{sec:dynamical-ybe-sl2}

\begin{definition}[Dynamical Yang--Baxter equation]
\label{def:dybe}
\index{dynamical Yang--Baxter equation|textbf}
The \emph{dynamical Yang--Baxter equation}
for $R(u, \lambda) \in \mathrm{End}(V \otimes V)$:
\begin{equation}\label{eq:dybe}
R_{12}(u, \lambda)\,
R_{13}(u{+}v, \lambda{-}h^{(2)})\,
R_{23}(v, \lambda)
=
R_{23}(v, \lambda{-}h^{(1)})\,
R_{13}(u{+}v, \lambda)\,
R_{12}(u, \lambda{-}h^{(3)}).
\end{equation}
\end{definition}

\begin{proposition}[DYBE reduces to Fay;
\ClaimStatusProvedHere]
\label{prop:dybe-reduces-to-fay}
\index{Fay trisecant identity!and dynamical YBE}
The dynamical Yang--Baxter equation for the
$\mathfrak{sl}_2$ elliptic $R$-matrix
\textup{(}Definition~\textup{\ref{def:elliptic-r})}
is equivalent to the Fay trisecant identity
\textup{(}Theorem~\textup{\ref{thm:fay})} plus the
quasi-periodicity of $\theta$.
\end{proposition}

\begin{proof}
On the weight-$1$ subspace of $V^{\otimes 3}$,
the nontrivial matrix entries of the
DYBE~\eqref{eq:dybe}, after clearing common
denominators, reduce to identities of the form
$\theta(a{-}b)\,\theta(c{-}d)
- \theta(a{-}c)\,\theta(b{-}d)
+ \theta(a{-}d)\,\theta(b{-}c) = 0$
for appropriate specializations $a, b, c, d$
depending on $u, v, \lambda, \eta$.  The
remaining components follow from the Weyl group
symmetry $R(u,\lambda)
= R(u,-\lambda)|_{+\leftrightarrow -}$
and the quasi-periodicity
$\theta(z{+}1|\tau) = -\theta(z|\tau)$.
\end{proof}

\begin{proposition}[DYBE encodes elliptic bar
nilpotency; \ClaimStatusProvedHere]
\label{prop:dybe-bar-nilpotency}
\index{bar complex!elliptic!nilpotency via DYBE}
The bar complex
$\B^{\mathrm{ell}}(E_{q,p}(\mathfrak{sl}_2))$
satisfies $d^2 = 0$ on $\overline{C}_3(E_\tau)$
if and only if $R(u,\lambda)$ satisfies the DYBE.
The dynamical shifts $\lambda \to \lambda - h^{(i)}$
arise from the $B$-cycle monodromy of the elliptic
propagator.
\end{proposition}


\subsection{Shuffle algebra presentation}
\label{subsec:shuffle-algebra}
\label{sec:shuffle-algebra}

\begin{definition}[Shuffle algebra]
\label{def:shuffle-algebra}
\index{shuffle algebra|textbf}
Let $\zeta(u) = \theta(u{+}\eta)/\theta(u)$.  The
\emph{shuffle algebra} $\mathcal{S}_\zeta$ has
$\mathcal{S}_n = \bC(z_1,\ldots,z_n)^{S_n}$
with product
\[
(f \star g)(z_1, \ldots, z_{m+n})
= \sum_{\sigma \in \mathrm{Sh}(m,n)}
\prod_{\substack{i \leq m \\ j > m}}
\zeta(z_{\sigma(i)} - z_{\sigma(j)})
\cdot f(z_{\sigma(1)}, \ldots, z_{\sigma(m)})
\cdot g(z_{\sigma(m+1)}, \ldots, z_{\sigma(m+n)}).
\]
The residue pairing gives $\mathcal{S}_n \cong
\B_n(\cA)^*$, identifying the shuffle product with
the deshuffle coproduct on $\B(\cA)$
(Corollary~\ref{cor:bar-is-dgcoalg}).
\end{definition}


\subsection{Comparison: rational, trigonometric, elliptic}
\label{subsec:rat-trig-ell}
\label{sec:rat-trig-ell-comparison}

\begin{table}[ht]
\centering
\caption{Comparison across three regimes}
\label{tab:rat-trig-ell}
\renewcommand{\arraystretch}{1.35}
{\footnotesize
\begin{tabular}{|l|c|c|c|}
\hline
& \textbf{Rational} & \textbf{Trigonometric}
& \textbf{Elliptic} \\ \hline
\textbf{Algebra}
& Yangian $Y(\fg)$
& Qtm.\ affine $U_q(\hat{\fg})$
& Toroidal
$U_{q,t}(\hat{\hat{\fg}})$ \\ \hline
\textbf{Curve}
& $\bC$
& $\bC^*$
& $E_\tau$ \\ \hline
\textbf{Propagator}
& $dz/z$
& $dz/(1 - z)$
& $d\log\theta_1(z|\tau)$ \\ \hline
\textbf{$R$-matrix}
& $1 - \hbar P/u$
& $(q{-}q^{-1})/(u{-}u^{-1})$
& $\theta(u{+}\eta)/\theta(\eta)$ \\ \hline
\textbf{Arnold id.}
& Partial fractions
& $q$-partial fractions
& Fay trisecant \\ \hline
\textbf{Koszul dual}
& $Y_{-\hbar}(\fg)$
& $U_{q^{-1}}(\hat{\fg})$
& $U_{q^{-1},t^{-1}}(\hat{\hat{\fg}})$
(conj.) \\ \hline
\textbf{Bar curv.}
& $m_0 = 0$
& $m_0 \propto \log q$
& $m_0 \propto E_2(\tau)$ \\ \hline
\textbf{Dyn.\ par.}
& absent
& absent
& $\lambda \in \mathfrak{h}^*$ \\ \hline
\end{tabular}
}
\end{table}

The degeneration chain is
\begin{equation}\label{eq:degeneration-chain}
\text{Elliptic}
\;\xrightarrow{\;\tau \to i\infty\;}
\text{Trigonometric}
\;\xrightarrow{\;q \to 1\;}
\text{Rational}.
\end{equation}
As $\tau \to i\infty$:
$\theta_1(z|\tau) \to 2\sin(\pi z)$ and the Eisenstein
corrections collapse to constants.  As $q \to 1$:
$dz/(1-z) \to dz/z$ and
$R_q(u) \to 1 - \hbar P/u$.
The bar complex degenerates accordingly:
$\B^{\mathrm{ell}} \to \B^{\mathrm{trig}}
\to \B^{\mathrm{rat}}$.


\subsection{Costello's M2-brane model}
\label{subsec:m2-brane}
\label{sec:m2-brane-double-loop}

\begin{definition}[5d NCCS theory {\cite{Costello17}};
\ClaimStatusProvedElsewhere]
\label{def:5d-nccs}
\index{noncommutative Chern--Simons!5d}
\index{M2-brane!double-loop algebra}
On $M_5 = \mathbb{R}_t \times \bC^2_{z_1,z_2}$ with
$\Omega$-background parameter $\varepsilon$, the 5d
noncommutative Chern--Simons theory has action
\begin{equation}\label{eq:5d-nccs-action}
S_{5d}(A)
= \frac{1}{2\delta}\int \mathrm{Tr}
\bigl(A \wedge dA\bigr)\,dz_1\,dz_2
+ \frac{1}{3\delta}\int \mathrm{Tr}
\bigl(A \wedge A *_\varepsilon A\bigr)\,dz_1\,dz_2,
\end{equation}
where $*_\varepsilon$ is the Moyal star product.  An
M2-brane wrapping
$\mathbb{R}_t \times \{(z_1,z_2) = (0,0)\}$ is an
instanton particle.
\end{definition}

The boundary algebra on the M2-brane is the double-current
algebra $\fg_{\mathrm{dbl}} = \mathfrak{gl}_r
\otimes \bC[u,v]$ with bracket
$[J^a_{m,n}, J^b_{p,q}]
= f^{ab}{}_c\,J^c_{m+p,n+q}$.
The two polynomial indices $(m,n)$ encode
the double loop: $u$ for $z_1$, $v$ for $z_2$.

\begin{proposition}[Classical Koszul duality for M2
boundary; \ClaimStatusProvedElsewhere]
\label{prop:m2-classical-koszul}
\index{Koszul duality!M2-brane}
$\Omega\B(U(\fg_{\mathrm{dbl}}))
\simeq C^\bullet(\fg_{\mathrm{dbl}})$,
the Chevalley--Eilenberg cochains.  The $L_\infty$
operations: $\ell_2(\epsilon^a_{m,n}, \epsilon^b_{p,q})
= f^{ab}{}_c\,\epsilon^c_{m+p,n+q}$
and $\ell_k = 0$ for $k \neq 2$
\textup{(}at the strict classical level\textup{)}.
In Costello's model, holographic duality \emph{is}
Koszul duality: the boundary algebra and the bulk
$L_\infty$ algebra are Koszul dual.  Bar-cobar inversion
$\Omega\B(\cA_\partial) \simeq \cA_\partial$ recovers
the boundary \textup{(}Theorem~B\textup{)};
the bulk is the derived centre
$C^\bullet_{\mathrm{ch}}(\cA_\partial, \cA_\partial)$
\textup{(}AP25, AP34\textup{)}.
\end{proposition}

\begin{conjecture}[Four-fold matching;
\ClaimStatusConjectured]
\label{prop:m2-four-fold-matching}
\index{M2-brane!sector/level/genus matching}
The Koszul duality $\cA_{\mathrm{bulk}}
= \Omega\B(\cA_{\mathrm{M2}})$ respects four
matchings: sector-wise \textup{(}adjoint labels\textup{)},
level-wise \textup{(}double-loop bidegree $(m,n)$
preserved\textup{)}, $N$-wise \textup{(}finite-$N$
quotient by trace relations\textup{)}, and genus-wise
\textup{(}bulk $g$-loop $= \hbar^{g-1}$ terms in
$L_\infty$ operations\textup{)}.
\end{conjecture}


%% ---------------------------------------------------------------
%% BRIDGE 1: E as shared geometry
%% ---------------------------------------------------------------

\subsection{The elliptic curve as shared geometry}
\label{subsec:bridge-shared-E}
\index{elliptic curve!shared geometry bridge}

\begin{remark}[Bridge: algebras on $E$ and algebras from $K3 \times E$]
\label{rem:bridge-shared-E}
The elliptic curve $E_\tau$ enters the remainder of this chapter in a second
capacity.  Acts~I and~II study chiral algebras \emph{on} $E_\tau$:
the propagator $d\log\theta_1(z|\tau)$, the Eisenstein corrections to
the bar differential, the toroidal algebra $U_{q,t}(\hat{\hat{\mathfrak{g}}})$
with $\tau = \log t / \log q$, and Felder's elliptic $R$-matrix
$R(u,\lambda)$ built from the same theta function.
The K3 chiral algebra $\cA_{K3}$ (Act~III) lives on a K3~surface, but
the Calabi--Yau threefold $K3 \times E$ (Act~IV) is fibered
\emph{over} $E$, and Costello--Li's holomorphic twist reduction along K3
produces a boundary chiral algebra $A_E$ \emph{on} $E_\tau$ with
$c(A_E) = \chi(K3) = 24$ free bosons from harmonic forms on K3
(Construction~\ref{constr:costello-li-k3xe}).

Thus: the toroidal and elliptic quantum group algebras of Act~II, and
the Costello--Li boundary algebra $A_E$ of Act~IV, are all chiral
algebras on the \emph{same} curve $E_\tau$.  They share the propagator
$d\log\theta_1$, the Fay trisecant relation governing $d^2 = 0$
(Proposition~\ref{prop:fay-implies-d-squared}), and the Eisenstein
filtration of the elliptic bar complex
(Theorem~\ref{thm:elliptic-vs-rational}).
The difference is the algebra: $A_E$ has $24$~generators
with $\kappa(A_E) = 24$ (rank of the free-boson lattice,
AP48) and shadow class~G (Gaussian, $r_{\max} = 2$), while
$E_{q,p}(\mathfrak{sl}_N)$ has $N^2 - 1$~generators with
dynamical parameter $\lambda \in \mathfrak{h}^*$.
\end{remark}


%% ===============================================================
%% ACT III: THE K3 CHIRAL ALGEBRA
%% ===============================================================

\section{The K3 chiral algebra}
\label{sec:k3-chiral-algebra}
\index{K3 surface!chiral algebra|textbf}
\index{N=4 superconformal algebra|textbf}

\subsection{The $\cN = 4$ superconformal algebra at
$c = 6$}
\label{subsec:n4-sca}

\begin{definition}[Small $\cN = 4$ SCA at $c = 6$]
\label{def:n4-sca-c6}
\index{N=4 superconformal algebra!c=6}
The \emph{small $\cN = 4$ superconformal algebra} at
$c = 6$ (equivalently, $\mathrm{SU}(2)_R$ level
$k_R = 1$, since $c = 6k_R$) has eight primary
generators:
\begin{center}
\small
\renewcommand{\arraystretch}{1.2}
\begin{tabular}{llcl}
\textbf{Generator} & \textbf{Weight $h$}
  & \textbf{Parity} & $J^3$ \textbf{charge} \\
\hline
$T$ & $2$ & bosonic & $0$ \\
$G^+, G^-$ & $3/2$ & fermionic & $\pm 1/2$ \\
$\tilde{G}^+, \tilde{G}^-$
  & $3/2$ & fermionic & $\mp 1/2$ \\
$J^{++}, J^{--}$ & $1$ & bosonic & $\pm 1$ \\
$J^3$ & $1$ & bosonic & $0$
\end{tabular}
\end{center}
The key OPEs at $c = 6$, $k_R = 1$:
\begin{align}
T(z)\,T(w) &\sim \frac{3}{(z-w)^4}
  + \frac{2T}{(z-w)^2}
  + \frac{\partial T}{z-w},
  \label{eq:n4-TT} \\
J^3(z)\,J^3(w)
  &\sim \frac{1/2}{(z-w)^2},
  \label{eq:n4-J3J3} \\
J^{++}(z)\,J^{--}(w)
  &\sim \frac{1}{(z-w)^2}
  + \frac{2J^3}{z-w},
  \label{eq:n4-JppJmm} \\
G^+(z)\,G^-(w)
  &\sim \frac{2}{(z-w)^3}
  + \frac{2J^3}{(z-w)^2}
  + \frac{T + \partial J^3}{z-w}.
  \label{eq:n4-GpGm}
\end{align}
\end{definition}


\subsection{Modular characteristic}
\label{subsec:kappa-k3}

\begin{proposition}[Modular characteristic of the K3
sigma model; \ClaimStatusProvedHere]
\label{prop:kappa-k3}
\index{modular characteristic!K3}
$\kappa(\cA_{K3}) = 2 = \dim_\bC(K3)$.
This is verified by five independent paths:
\begin{enumerate}[label=\textup{(\roman*)}]
\item \emph{Geometric.}
  $\kappa(\Omega^{\mathrm{ch}}(\mathrm{CY}_d)) = d$
  by the Chern character computation.
\item \emph{Character.}
  $F_1 = \kappa/24 = 1/12$.
\item \emph{Complementarity.}
  $\kappa + \kappa' = 0$ for CY sigma models
  \textup{(}AP24: this holds for free-field and
  KM-type algebras; for W-algebras
  $\kappa + \kappa' = \rho K$\textup{)},
  giving $\kappa' = -2$.
\item \emph{Kummer.}
  At the orbifold point $K3 = T^4/\bZ_2$: the
  orbifold preserves $\kappa = 2$.
\item \emph{Additivity.}
  $\kappa(K3 \times E) = \kappa(K3) + \kappa(E)
  = 2 + 1 = 3 = \dim_\bC(K3 \times E)$
  \textup{(}Corollary~\textup{\ref{cor:shadow-extraction})}.
\end{enumerate}
\end{proposition}

\begin{proof}
The index-theoretic computation
$\kappa(\Omega^{\mathrm{ch}}(\mathrm{CY}_d)) = d$
reduces to $\chi(K3)/12 = 24/12 = 2$ via the Chern
character of the chiral de~Rham complex.  Curve
independence fixes $\kappa$ across K3 moduli.
Additivity is
Proposition~\ref{prop:independent-sum-factorization}.
\end{proof}

\begin{remark}[AP48: $\kappa(K3) = 2 \neq c/2 = 3$]
\label{rem:ap48-k3}
\index{AP48!K3 example}
The $\cN = 4$ Ward identities constrain $\kappa$ below
the Virasoro formula $c/2 = 3$.  The Gepner model
$(2)^4$ gives $\kappa_{\mathrm{Gepner}} = 4 \times 3/4
= 3 \neq 2$: a different algebra from the $\cN = 4$ SCA
(AP48: $\kappa$ depends on the full algebra, not the
Virasoro subalgebra).
\end{remark}


\subsection{Bar complex and Koszulness}
\label{subsec:bar-k3}

\begin{proposition}[Bar complex of the $\cN = 4$ SCA;
\ClaimStatusProvedHere]
\label{prop:bar-n4-sca}
\index{bar complex!N=4 SCA}
$B(\cA_{K3})$ has: $\dim B^1 = 8$
\textup{(}one desuspended generator per primary;
AP45: $|s^{-1}v| = |v| - 1$\textup{)};
$d\colon B^2 \to B^1$ is rank~$16$
\textup{(}AP19: $r$-matrix poles one below OPE
poles\textup{)};
and the algebra is chirally Koszul
\textup{(}Proposition~\textup{\ref{prop:pbw-universality}}:
freely strongly generated, no null vectors in the
universal vacuum module\textup{)}.
Shadow depth: class~M ($r_{\max} = \infty$),
because the critical discriminant
$\Delta = 8\kappa\,S_4 \neq 0$.
\end{proposition}


\subsection{Koszul dual}
\label{subsec:koszul-dual-k3}

\begin{proposition}[Koszul dual of $\cA_{K3}$;
\ClaimStatusProvedHere]
\label{prop:koszul-dual-k3}
\index{Koszul dual!K3 chiral algebra}
$\kappa(\cA_{K3}^!) = -2$
\textup{(}complementarity: $\kappa + \kappa' = 0$
for CY sigma models \textup{(}AP24: not universal;
for W-algebras $\kappa + \kappa' = \rho K$\textup{))}.
$c(\cA_{K3}^!) = -6$
\textup{(}$k_R' = -1$, so $c' = 6k_R' = -6$\textup{)}.
$F_g(\cA^!) = -2 \cdot \lambda_g^{\mathrm{FP}}$
at genus~$1$ unconditionally; at $g \geq 2$ on
the uniform-weight lane \textup{(}AP32\textup{)}.
Verdier intertwining \textup{(}Theorem~A\textup{)}:
$D_{\mathrm{Ran}}(B(\cA_{K3})) \simeq B(\cA_{K3}^!)$.
\end{proposition}


\subsection{Elliptic genus and Mathieu moonshine}
\label{subsec:elliptic-genus-k3}
\label{subsec:k3-elliptic-genus}
\label{sec:mock-modular-bkm}

\begin{definition}[K3 elliptic genus]
\label{def:k3-elliptic-genus}
\index{elliptic genus!K3|textbf}
$Z_{K3}(\tau,z) = 2\,\phi_{0,1}(\tau,z)$,
where $\phi_{0,1}$ is the unique weak Jacobi form of
weight~$0$, index~$1$ (Eichler--Zagier normalization:
$\phi_{0,1}(\tau,0) = 12$, AP38).
At $z = 0$: $Z_{K3}(\tau,0) = 24 = \chi(K3)$.
At $q^0$: $2y + 20 + 2y^{-1}$ (the $\chi_y$-genus).
\end{definition}

\begin{theorem}[Mathieu moonshine;
\ClaimStatusProvedElsewhere]
\label{thm:mathieu-moonshine}
\index{Mathieu moonshine|textbf}
The $\cN = 4$ decomposition
$Z_{K3} = 20 \cdot \mathrm{ch}_{\mathrm{massless}}
+ \sum_{n \geq 1} A_n \cdot
\mathrm{ch}_{\mathrm{massive},n}$
has coefficients $A_n$ that are dimensions of
$M_{24}$-representations
\textup{(}Eguchi--Ooguri--Tachikawa 2010,
proved Gannon 2016\textup{)}:
\begin{center}
\small
\renewcommand{\arraystretch}{1.2}
\begin{tabular}{ccccccc}
$n$ & $1$ & $2$ & $3$ & $4$ & $5$ & $6$ \\ \hline
$A_n$ & $90$ & $462$ & $1540$ & $4554$
  & $11592$ & $27830$
\end{tabular}
\end{center}
The overall factor $2 = \kappa(\cA_{K3})$ is the
modular characteristic:%
\label{rem:factor-2-is-kappa}%
\begin{equation}\label{eq:moonshine-multiplier}
A_n = \kappa(\cA_{K3}) \cdot \dim(\rho_n),
\end{equation}
where $\rho_n \in \mathrm{Irr}(M_{24})$ with
$\dim(\rho_1) = 45$, $\dim(\rho_2) = 231$,
$\dim(\rho_3) = 770$.
The bar complex detects $\kappa = 2$ at arity~$2$;
the $M_{24}$ action is an internal symmetry of the
massive BPS sector, orthogonal to the bar construction.
\end{theorem}

\begin{remark}[Mock modular form from K3]
\label{rem:mock-modular-k3}
\index{mock modular form!K3}
The massive multiplicities assemble into the mock
modular form
$h(\tau) = 2q^{-1/8}(-1 + 45q + 231q^2 + 770q^3
+ \cdots)$,
with shadow $S(\tau) = 24\,\eta(\tau)^3$ (weight~$3/2$
cusp form).  The polar coefficient $-2 = -\kappa$
reappears as the modular characteristic.
\end{remark}

\begin{remark}[Lattice VOA, Gepner, and chiral de~Rham]
\label{rem:lattice-voa-k3}
\index{lattice VOA!K3}
\index{chiral de Rham complex!K3}
\index{Gepner model}
The K3 lattice $L_{K3} = H^2(K3,\bZ) \cong U^3 \oplus
(-E_8)^2$ has rank~$22$ and signature $(3,19)$.  The
negative-definite sublattice $(-E_8)^2$ gives a lattice
VOA with $c = 16$, $\kappa = 16$ (AP48:
$\kappa = \mathrm{rank}$ for lattice VOAs).
The chiral de~Rham complex $\Omega^{\mathrm{ch}}(K3)$
(Malikov--Schechtman--Vaintrob) has $c = 0$ and
recovers $Z_{K3} = 2\,\phi_{0,1}$ via
the Borisov--Libgober formula.
The Gepner model $(2)^4$ for the Fermat quartic
has chiral ring $\mathrm{Jac}(W)$ with
$\dim = 81$ before orbifold.
\end{remark}


%% ===============================================================
%% ACT IV: K3 x E AS CALABI--YAU THREEFOLD
%% ===============================================================

\section{The Calabi--Yau threefold
\texorpdfstring{$K3 \times E$}{K3 x E}}
\label{sec:k3xe}
\label{sec:k3xe-geometry}
\index{K3 x E|textbf}

\subsection{Geometry}
\label{subsec:geometry-k3xe}

\begin{proposition}[Hodge diamond;
\ClaimStatusProvedElsewhere]
\label{prop:hodge-k3xe}
\index{Hodge diamond!K3 x E}
The K\"unneth formula gives:
\[
\renewcommand{\arraystretch}{1.1}
\begin{array}{ccccccc}
& & & 1 & & & \\
& & 1 & & 1 & & \\
& 1 & & 21 & & 1 & \\
1 & & 21 & & 21 & & 1 \\
& 1 & & 21 & & 1 & \\
& & 1 & & 1 & & \\
& & & 1 & & &
\end{array}
\]
$h^{1,1} = h^{2,1} = 21$.
$\chi(K3 \times E) = 24 \cdot 0 = 0$.
$\mathrm{rank}\,K_0 = 48$.
$\dim \mathrm{HH}^* = 96$.
\end{proposition}


\subsection{Shadow obstruction tower}
\label{subsec:shadow-k3xe}

\begin{proposition}[Shadow amplitudes of $K3 \times E$;
\ClaimStatusProvedHere]
\label{prop:shadow-k3e}
\index{shadow amplitudes!K3 x E}
$\kappa(\Omega^{\mathrm{ch}}(K3 \times E))
= \dim_\bC(K3 \times E) = 3$.
\textup{(}AP48: this coincides with $c/2 = 3$ only
because $\dim_\bC = 3$.\textup{)}
The elliptic genus $Z_{K3 \times E} = Z_{K3} \cdot Z_E
= 0$ \textup{(}since $\chi_y(E) = 0$\textup{)},
but $\kappa \neq 0$:
the bar complex curvature detects geometry invisible
to the elliptic genus \textup{(}AP31\textup{)}.

The shadow amplitudes
\textup{(}Theorem~D, in the convention
$\sum_{g \geq 1} F_g\,\hbar^{2g}$, AP22\textup{)}:
\begin{equation}\label{eq:shadow-k3e}
F_g(K3 \times E) = 3 \cdot
\frac{(2^{2g-1}-1)\,|B_{2g}|}{2^{2g-1}\,(2g)!}.
\end{equation}
\begin{center}
\small
\renewcommand{\arraystretch}{1.2}
\begin{tabular}{cccccc}
$g$ & $1$ & $2$ & $3$ & $4$ & $5$ \\ \hline
$F_g$ & $\frac{1}{8}$ & $\frac{7}{1920}$
  & $\frac{31}{322560}$
  & $\frac{127}{51609600}$
  & $\frac{73}{1167851520}$
\end{tabular}
\end{center}
\end{proposition}

\begin{proof}
$\kappa = 3$ by the Chern character computation for
$\Omega^{\mathrm{ch}}(\mathrm{CY}_d)$ with $d = 3$,
or by additivity $\kappa(K3) + \kappa(E) = 2 + 1 = 3$.
Each $F_g$ is verified by the Bernoulli formula, the
$\hat{A}$-generating function
$(t/2)/\sin(t/2) - 1 = \sum \lambda_g^{\mathrm{FP}}
\,t^{2g}$, and the additivity of $\kappa$.
\end{proof}

\begin{proposition}[Convergence; \ClaimStatusProvedHere]
\label{prop:shadow-gf-convergence}
\index{shadow generating function!convergence}
$Z^{\mathrm{sh}}(t) := 3\bigl((t/2)/\sin(t/2) - 1\bigr)
= \sum_{g \geq 1} F_g\,t^{2g}$ has convergence radius
$R = 2\pi$ \textup{(}poles of $1/\sin(t/2)$
at $t = \pm 2\pi$\textup{)}.  The Borel transform is
entire: no resurgence, no Stokes phenomenon.  This
contrasts with the topological string, which is
Gevrey-$1$ divergent: $F_g^{\mathrm{top}} \sim (2g)!$.
\end{proposition}


\subsection{The four $\kappa$ invariants}
\label{subsec:four-kappas}

\begin{remark}[Four distinct $\kappa$ values]
\label{rem:four-kappas}
\index{modular characteristic!four variants}
\begin{center}
\small
\renewcommand{\arraystretch}{1.2}
\begin{tabular}{lll}
\textbf{Symbol} & \textbf{Value}
  & \textbf{Source} \\ \hline
$\kappa(\cA)$ & $3$
  & Modular characteristic (Theorem~D) \\
$\kappa(K3)$ & $2$
  & K3 sigma model
    (Proposition~\ref{prop:kappa-k3}) \\
$\kappa(A_E)$ & $24$
  & Costello--Li boundary algebra
    (Constr.~\ref{constr:costello-li-k3xe}) \\
$\kappa_{\mathrm{BKM}}$ & $5$
  & $\operatorname{wt}(\Phi_{10})/2$
\end{tabular}
\end{center}
These coincide for Virasoro ($\kappa = c/2$ in all
r\^oles) but diverge for $K3 \times E$ (AP20, AP48).
The intrinsic invariant is $\kappa(\cA) = 3$.
\end{remark}


\subsection{Enumerative geometry}
\label{subsec:enum-geom}
\label{sec:enum-geom-k3xe}
\label{sec:bps-bh-k3xe}

The Yau--Zaslow formula (proved Pandharipande--Thomas
2016) gives genus-$0$ BPS invariants:
\begin{equation}\label{eq:yau-zaslow}
\sum_{h \geq 0} n^0_h\,q^h
= \prod_{n \geq 1} \frac{1}{(1 - q^n)^{24}}.
\end{equation}
Selected values:
\begin{center}
\small
\renewcommand{\arraystretch}{1.2}
\begin{tabular}{lccccccc}
$h$ & $0$ & $1$ & $2$ & $3$ & $4$ & $5$ & $6$ \\
\hline
$n^0_h$ & $1$ & $24$ & $324$ & $3200$ & $25650$
  & $176256$ & $1073720$ \\
$n^1_h$ & $0$ & $-2$ & $-54$ & $-800$ & $-8550$
  & $-73440$ & $-536860$
\end{tabular}
\end{center}
Since $\chi(K3 \times E) = 0$: the DT/PT correspondence
is clean ($M(-q)^0 = 1$), the BCOV anomaly coefficient
$\chi/24 = 0$, and the MacMahon prefactor is trivial.

\subsection{Twisted holography}
\label{subsec:twisted-holography}
\label{sec:twisted-holography-k3xe}

\begin{construction}[Costello--Li KS boundary algebra]
\label{constr:costello-li-k3xe}
\label{prop:boundary-sigma-ratio}%
\label{rem:leech-connection}%
\index{Kodaira--Spencer theory!K3 x E}
\index{twisted holography!K3 x E}
The HT twist of Type~IIB on $K3 \times E$, reduced
along K3, produces a boundary chiral algebra $A_E$ on $E$
with $c(A_E) = \chi(K3) = 24$
($24$ free bosons from harmonic forms on K3) and
$\kappa(A_E) = 24$ (AP48: rank of the free-boson lattice).
The holographic modular Koszul datum is
$H(K3 \times E) = (A_E, A_E^!, C, r(z), \Theta_{A_E},
\nabla^{\mathrm{hol}})$
with $r(z) = \Omega/z$ ($24$-dimensional Casimir)
and shadow class~G.
\end{construction}


%% ---------------------------------------------------------------
%% BRIDGE 2: Fourier--Jacobi as genus escalation
%% ---------------------------------------------------------------

\subsection{Fourier--Jacobi as genus escalation}
\label{subsec:bridge-fourier-jacobi}
\index{Fourier--Jacobi expansion!genus escalation bridge}

\begin{proposition}[Bridge: elliptic $R$-matrix data determines BKM
root system; \ClaimStatusProvedElsewhere]
\label{prop:bridge-fj-genus-escalation}
The Fourier--Jacobi expansion of the Igusa cusp form
$\Phi_{10}(\Omega)
= \sum_{m \geq 1} \phi_m(\tau,z)\,p^m$
uses Jacobi forms $\phi_m(\tau,z)$ on $E_\tau$.  The leading
coefficient
\begin{equation}\label{eq:bridge-phi-1}
\phi_1(\tau,z) = \eta(\tau)^{18}\,\vartheta_1(\tau,z)^2
\end{equation}
is built from the same theta function $\vartheta_1(\tau,z)$ that defines the
elliptic propagator $d\log\vartheta_1(z|\tau)$
\textup{(}Definition~\textup{\ref{def:elliptic-propagator})} and
Felder's elliptic $R$-matrix entries $a(u) = \theta(u{+}\eta)/\theta(\eta)$
\textup{(}Definition~\textup{\ref{def:elliptic-r})}.

The passage from genus-$1$ data \textup{(}Jacobi forms on $E_\tau$\textup{)}
to genus-$2$ data \textup{(}the Siegel form $\Phi_{10}$ on $\mathfrak{H}_2$\textup{)}
is the Borcherds multiplicative lift.  Its input is the K3 elliptic genus
$Z_{K3} = 2\,\phi_{0,1}$
\textup{(}Definition~\textup{\ref{def:k3-elliptic-genus})}, and the exponents
$c_0(4nm - \ell^2)$ in the infinite product~\eqref{eq:borcherds-lift} are precisely
the root multiplicities of the BKM superalgebra $\mathfrak{g}_{\Phi_{10}}$
\textup{(}Proposition~\textup{\ref{prop:bkm-roots})}.

In the factorization $\eta^{18} = \eta^{24} \cdot \eta^{-6}$:
$\eta^{24} = \Delta$ is the Ramanujan discriminant
\textup{(}the $24$~free bosons of $A_E$, with $\kappa(A_E) = 24$\textup{)}, and
$\eta^{-6} = \eta^{-2\kappa}$ with $\kappa = \kappa(\Omega^{\mathrm{ch}}(K3 \times E))
= 3$.  The shadow obstruction tower detects $\kappa = 3$
\textup{(}Proposition~\textup{\ref{prop:shadow-k3e})};
the Borcherds lift escalates this genus-$1$ datum to the genus-$2$ Siegel modular form.
\end{proposition}


%% ===============================================================
%% ACT V: THE BKM CROWN
%% ===============================================================

\section{The BKM crown}
\label{sec:bkm-k3e}
\index{Borcherds--Kac--Moody algebra|textbf}
\index{Igusa cusp form|textbf}

\subsection{The Igusa cusp form}
\label{subsec:igusa}

\begin{definition}[Igusa cusp form]
\label{def:igusa-cusp-form}
\index{Igusa cusp form!definition}
$\Phi_{10}$ is the unique Siegel cusp form of weight~$10$
on $\operatorname{Sp}(4,\bZ)$:
$\dim S_{10}(\operatorname{Sp}(4,\bZ)) = 1$.
Three constructions:
\begin{enumerate}[label=\textup{(\roman*)}]
\item \emph{Even theta characteristics.}
$\Phi_{10}(\Omega) = -2^{-12}
\prod_{\text{even}\,m} \vartheta[m](\Omega)$,
the product over the $10$ even theta
characteristics of genus~$2$.
\item \emph{Fourier--Jacobi expansion.}
$\Phi_{10}(\Omega) = \sum_{m \geq 1}
\phi_m(\tau,z)\,p^m$
with $p = e^{2\pi i \sigma}$,
$\Omega = \bigl(\begin{smallmatrix}
\tau & z \\ z & \sigma
\end{smallmatrix}\bigr)
\in \mathfrak{H}_2$, and
\begin{equation}\label{eq:phi-10-1}
\phi_1 = \phi_{10,1}(\tau,z)
= \eta(\tau)^{18}\,\vartheta_1(\tau,z)^2.
\end{equation}
\item \emph{Borcherds lift.}
\begin{equation}\label{eq:borcherds-lift}
\Phi_{10}(\Omega)
= p\,q\,y \prod_{\substack{(n,\ell,m) > 0}}
\bigl(1 - q^n\,y^\ell\,p^m\bigr)^{c_0(4nm - \ell^2)},
\end{equation}
where $c_0(D)$ are the discriminant-indexed Fourier
coefficients of $Z_{K3} = 2\,\phi_{0,1}$
\textup{(}AP38: Eichler--Zagier convention\textup{)}.
\end{enumerate}
\end{definition}

The DVV formula (Dijkgraaf--Verlinde--Verlinde):
\begin{equation}\label{eq:dvv}
Z_{\frac{1}{4}\text{-BPS}}(\Omega)
= \frac{1}{\Phi_{10}(\Omega)}
= \sum_{N \geq 0}\,p^N\,
\chi\bigl(\operatorname{Sym}^N(K3); \tau, z\bigr).
\end{equation}
The $\frac{1}{4}$-BPS dyon with discriminant $\Delta$
has $S_{\mathrm{BH}} = 4\pi\sqrt{\Delta}$.


\subsection{The Borcherds lift: genus-$1$ to genus-$2$}
\label{subsec:borcherds-lift}

\begin{theorem}[Borcherds lift;
\ClaimStatusProvedElsewhere]
\label{thm:borcherds-lift-k3e}
\index{Borcherds lift|textbf}
The Borcherds multiplicative lift sends the genus-$1$
datum $Z_{K3} = 2\,\phi_{0,1}$ to the genus-$2$ Siegel
form $\Phi_{10}$ via the infinite
product~\eqref{eq:borcherds-lift}.  The exponent
$c_0(D)$ at discriminant $D$ is the root multiplicity of
the BKM superalgebra $\mathfrak{g}_{\Phi_{10}}$.
The factorization $\eta^{18} = \eta^{24} \cdot \eta^{-6}$
separates the Ramanujan discriminant from
$\eta^{-2\kappa}$ with $\kappa = 3$.
\end{theorem}


\subsection{BKM root multiplicities}
\label{subsec:bkm-roots}

\begin{proposition}[Root multiplicities;
\ClaimStatusProvedElsewhere]
\label{prop:bkm-roots}
\index{BKM algebra!root multiplicities}
$\mathfrak{g}_{\Phi_{10}}$ has root multiplicities from
$Z_{K3}$:
\begin{itemize}
\item Real roots ($D = -1$, $\alpha^2 = 2$):
  $\mathrm{mult} = c_0(-1) = 2$.
\item Lightlike roots ($D = 0$, $\alpha^2 = 0$):
  $\mathrm{mult} = c_0(0) = 20$
  (from $2\phi_{0,1}$).
\item Imaginary roots ($D > 0$): unbounded.
  $c_0(3) = -64$, $c_0(4) = 108$,
  $c_0(7) = -513$, $c_0(8) = 808$.
\end{itemize}
Negative multiplicities at odd $D$ indicate fermionic
directions.
\end{proposition}


%% ---------------------------------------------------------------
%% BRIDGE 3: Hilbert schemes
%% ---------------------------------------------------------------

\subsection{The Hilbert scheme bridge}
\label{subsec:bridge-hilbert-schemes}
\index{Hilbert scheme!bridge to toroidal algebras}

\begin{theorem}[Toroidal action on $K$-theory of Hilbert schemes
{\cite{SV13}}; \ClaimStatusProvedElsewhere]
\label{thm:bridge-sv-toroidal}
\index{Schiffmann--Vasserot theorem}
\index{toroidal algebra!Hilbert scheme action}
Let $S$ be a smooth projective surface.  The toroidal algebra
$U_{q,t}(\hat{\hat{\mathfrak{gl}}}_1)$ acts on
$\bigoplus_{n \geq 0} K({\operatorname{Hilb}}^n(S))$
by correspondences on the nested Hilbert scheme
$\operatorname{Hilb}^{n,n+1}(S)$ \textup{(}Schiffmann--Vasserot\textup{)}.
Specializations:
\begin{enumerate}[label=\textup{(\roman*)}]
\item \emph{Nakajima operators} \textup{(}$q = 1$\textup{)}: the
  Heisenberg algebra $\bigoplus_{n \in \bZ} \bQ\,\mathfrak{a}_n$
  acts on $\bigoplus_n H^*(\operatorname{Hilb}^n(S))$ via
  Nakajima correspondences,
  with $\mathfrak{a}_{-n}|0\rangle$
  the class of the punctual Hilbert scheme.
\item \emph{Haiman's theorem} \textup{(}$S = \bC^2$\textup{)}:
  the DAHA $\ddot{H}_n(q,t)$, a quotient of
  $U_{q,t}(\hat{\hat{\mathfrak{gl}}}_1)$, acts on
  $K(\operatorname{Hilb}^n(\bC^2))$ with fixed-point basis
  indexed by partitions $\lambda \vdash n$, and the
  character of the natural representation is the Macdonald
  polynomial $\widetilde{H}_\lambda(x;q,t)$.
\item \emph{$S = K3$}:
  $\chi(\operatorname{Hilb}^n(K3)) = p_{24}(n)$
  \textup{(}$24$-coloured partitions\textup{)}, so
  $\sum_{n \geq 0} \chi(\operatorname{Hilb}^n(K3))\,p^n
  = \prod_{m \geq 1} (1 - p^m)^{-24}$.
\end{enumerate}
\end{theorem}

\begin{proposition}[DMVV from Hilbert scheme characters;
\ClaimStatusProvedElsewhere]
\label{prop:bridge-dmvv-hilb}
\index{DMVV formula!Hilbert scheme}
The DVV formula~\eqref{eq:dvv}
$1/\Phi_{10} = \sum_N p^N\,\chi(\operatorname{Sym}^N(K3);\tau,z)$
is the generating function of equivariant Euler characteristics of
$\operatorname{Hilb}^N(K3)$ refined by the $\mathrm{SU}(2)_R$~fugacity
$z$ and the $L_0$~fugacity $q = e^{2\pi i\tau}$.  The $24$~coloured
partitions arise from Nakajima's Heisenberg action
\textup{(}$24 = \chi(K3) = b_0 + b_2 + b_4 = 1 + 22 + 1$\textup{)}:
each colour corresponds to a cohomology class of K3, matching the
$24$~free bosons of the boundary algebra $A_E$
\textup{(}Construction~\textup{\ref{constr:costello-li-k3xe})}.
\end{proposition}

\begin{remark}[Synthesis: toroidal algebras meet BPS counting]
\label{rem:bridge-synthesis}
\index{Hilbert scheme!synthesis}
The three strands of this chapter converge on the Hilbert scheme.
The toroidal algebra $U_{q,t}(\hat{\hat{\mathfrak{g}}})$ of
\S\ref{subsec:toroidal-algebra} provides the symmetry algebra
acting on $K(\operatorname{Hilb}^n(S))$.
The elliptic $R$-matrix of
\S\ref{subsec:elliptic-quantum-groups} appears as the transition
matrix between stable-envelope bases of $K(\operatorname{Hilb}^n(S))$
(Aganagic--Okounkov).
The BKM root multiplicities $c_0(D)$ of
Proposition~\ref{prop:bkm-roots} are the Fourier coefficients of the
K3~elliptic genus $Z_{K3} = 2\,\phi_{0,1}$, which in turn
counts representations of the Heisenberg subalgebra on
$\bigoplus_n H^*(\operatorname{Hilb}^n(K3))$.
The shadow invariant $\kappa = 3$ measures the single-copy
chiral algebra $\Omega^{\mathrm{ch}}(K3 \times E)$; the BKM algebra
$\mathfrak{g}_{\Phi_{10}}$ encodes the second-quantized
(symmetric-product) extension, and the gap between them is the content
of the next subsection.
\end{remark}


\subsection{The second-quantization gap}
\label{subsec:second-quantization}

\begin{proposition}[BPS modular characteristic;
\ClaimStatusProvedHere]
\label{prop:kappa-bps-decomposition}
\index{second quantization!kappa decomposition}
\begin{equation}\label{eq:kappa-bps-decomposition}
\kappa_{\mathrm{BPS}}
= \kappa(K3) + \kappa(K3 \times E)
= 2 + 3 = 5
= \tfrac{1}{2}\operatorname{wt}(\Phi_{10}).
\end{equation}
The ratio
$\kappa_{\mathrm{BPS}}/\kappa_{\mathrm{ch}} = 5/3$
reflects the passage from single-copy to symmetric
product, specific to $K3 \times E$.
\end{proposition}


\subsection{The shadow--Siegel gap}
\label{subsec:shadow-siegel-gap}

\begin{theorem}[Shadow--Siegel gap;
\ClaimStatusProvedHere]
\label{thm:shadow-siegel-gap}
\index{shadow--Siegel gap|textbf}
The shadow obstruction tower does not produce
$\Phi_{10}$.  Four independent obstructions:
\begin{enumerate}[label=\textup{(O\arabic*)}]
\item \textbf{Categorical.}
$F_g \in \bQ$ is a tautological intersection number;
$\Phi_{10} \in \cO(\mathfrak{H}_2)$ is a Siegel
cusp form.

\item \textbf{Modular characteristic.}
The shadow uses $\kappa_{\mathrm{ch}} = 3$;
$\Phi_{10}$ has weight $10 = 2\kappa_{\mathrm{BPS}}$
with $\kappa_{\mathrm{BPS}} = 5$.

\item \textbf{Second quantization.}
$F_g$ is first-quantized (single copy of $\cA$).
$1/\Phi_{10}$ is second-quantized (DMVV symmetric
product over all $\operatorname{Sym}^N(K3)$).

\item \textbf{Schottky at $g \geq 4$.}
The Jacobian locus $\cJ_g$ has codimension
$(g-2)(g-3)/2$ in $\cA_g$ for $g \geq 4$.
\begin{center}
\small
\renewcommand{\arraystretch}{1.1}
\begin{tabular}{cccccc}
$g$ & $1$ & $2$ & $3$ & $4$ & $5$ \\ \hline
$\mathrm{codim}$ & $0$ & $0$ & $0$ & $1$ & $3$
\end{tabular}
\end{center}
For $g \leq 3$: the shadow sees all Siegel data.
At $g = 4$: the Schottky--Igusa form detects the
Jacobian locus, invisible to the tautological shadow.
\end{enumerate}
\end{theorem}

\begin{proof}
(O1): $F_g \in \bQ$ by definition; $\Phi_{10}$ is
holomorphic on $\mathfrak{H}_2$.
(O2): $\kappa_{\mathrm{ch}} = 3$ by
Proposition~\ref{prop:shadow-k3e};
$\operatorname{wt}(\Phi_{10}) = 10$ by
Definition~\ref{def:igusa-cusp-form}.
(O3): the DMVV formula~\eqref{eq:dvv}.
(O4): the codimension formula
$\mathrm{codim}(\cJ_g, \cA_g) = (g-2)(g-3)/2$
(Riemann, Schottky).
\end{proof}

\begin{remark}[Genus-$2$ Torelli bridge]
\label{rem:torelli-bridge}
\index{Torelli map!genus-2}
At genus~$2$, the Torelli map is birational: $\Phi_{10}$
restricts to $\cM_2$, and the shadow amplitude
$F_2 = 7/1920$ is its topological residue.
\end{remark}


%% ===============================================================
%% ACT VI: THE FRONTIER
%% ===============================================================

\section{Research programmes}
\label{sec:cy-research-programmes}
\label{sec:grand-synthesis-k3xe}
\label{subsec:cy-research-programmes}
\index{K3 x E!research programmes}

\begin{openproblem}[Bar complex detection of $M_{24}$]
\label{op:bar-m24}
\index{Mathieu moonshine!bar complex}
The bar complex detects $\kappa = 2$ at arity~$2$ and
the cubic shadow at arity~$3$.  Does the full shadow
obstruction tower $\Theta_{\cA_{K3}}$ encode the
$M_{24}$ action on the massive BPS sector?
The mock modular shadow $24\,\eta^3$ and the
bar-complex shadow connection share the feature of
encoding modular non-invariance; whether they are
related by a precise functor remains open.
\end{openproblem}

\begin{conjecture}[Universal moonshine multiplier;
\ClaimStatusConjectured]
\label{conj:universal-moonshine-multiplier}
\label{op:programme-b-moonshine}
\index{moonshine!universal multiplier}
For a holomorphic VOA~$V$ with finite automorphism
group~$G$:
$Z_g(\tau,z) = \kappa(V) \cdot
\sum_{n \geq 0} \dim(\rho_n)\,\mathrm{ch}_n(g)$
for $g \in G$.  For $V = \cA_{K3}$:
$\kappa = 2$, $A_n = 2 \cdot \dim(\rho_n)$.
\end{conjecture}

\begin{openproblem}[Second-quantization bridge]
\label{op:programme-c-second-quantization}
\label{conj:kappa-bps-universality}
\index{second quantization!shadow bridge}
The ratio $\kappa_{\mathrm{BPS}}/\kappa_{\mathrm{ch}}
= 5/3$ encodes the passage from first to second
quantization.  The DMVV formula
$1/\Phi_{10} = \mathrm{PE}[2\,\phi_{0,1}]$ expresses
the second-quantized partition function as the
plethystic exponential of the single-copy elliptic
genus.  What is the shadow tower of
$\mathrm{PE}[\cA]$ in terms of the shadow tower
of~$\cA$?
\end{openproblem}

\begin{openproblem}[Schottky shadow programme]
\label{op:programme-d-schottky}
\label{conj:shadow-taut-projection}
\index{Schottky problem!shadow tower}
The shadow tower lives on $\overline{\cM}_g$;
Siegel forms live on $\cA_g$.  For $g \leq 3$, the
Torelli map identifies the two.  For $g \geq 4$, the
Schottky--Igusa form (weight $8$ at $g = 4$, vanishing
on $\cJ_4$) detects the Jacobian locus, invisible to
the tautological shadow.
\end{openproblem}

\begin{openproblem}[Mock modularity and the bar complex]
\label{op:programme-e-mock-modularity}
\label{conj:mock-shadow-tower}
\index{mock modular form!bar complex}
The mock modular form $h(\tau)$ has shadow
$S(\tau) = 24\,\eta^3$ (weight $3/2$).  The bar-complex
shadow connection produces integer-weight objects.
Is there a metaplectic bar complex whose shadow
connection produces weight-$3/2$ objects?
\end{openproblem}

\begin{openproblem}[Factorization envelope of the K3
chiral algebra]
\label{op:fact-envelope-k3}
\label{op:programme-f-factorization-envelope}
\label{conj:modular-fact-envelope}
\index{factorization envelope!K3}
$U^{\mathrm{mod}}_X(\cA_{K3})$ (Nishinaka 2025/26,
Vicedo 2025) is not yet constructed.  The six-fold
platonic package $\Pi_X(\cA_{K3})$ is well-defined
(Construction~\ref{constr:platonic-package}); the
modular factorization envelope as the left adjoint of
$\mathrm{Prim}^{\mathrm{mod}}$ remains open.
\end{openproblem}

\begin{openproblem}[BKM algebras and scattering diagrams]
\label{op:programme-g-bkm-scattering}
\label{conj:scattering-tropical-shadow}
\index{scattering diagram!BKM}
The shadow obstruction tower $\Theta_{\cA}$ and the
KS automorphism $\mathfrak{S}$ are MC elements in
different Lie algebras.  The identification
``scattering $=$ shadow'' holds at the motivic Hall
algebra level (AP42: the naive BCH pair-commutator
does \emph{not} reproduce $\phi_{0,1}$ multiplicities).
The tropical limit of the scattering diagram should
connect Mok's log FM tropicalization
(Theorem~\ref{thm:planted-forest-tropicalization})
to BPS walls.
\end{openproblem}

\begin{openproblem}[Higher-dimensional CY shadows]
\label{op:programme-i-higher-dim}
\label{conj:cy-product-shadow}
\index{Calabi--Yau!higher-dimensional shadows}
By additivity:
$\kappa(K3 \times K3) = 4$,
$\kappa(K3 \times K3 \times E) = 5$.
The shadow class of a product is
$\max(\text{classes of factors})$ under $G < L < C < M$.
Is there a genus-$3$ Siegel form analogue of
$\Phi_{10}$ for $\mathrm{CY}_4$?
\end{openproblem}

\begin{openproblem}[Convergence versus divergence]
\label{op:programme-j-convergence}
\index{shadow partition function!convergence}
The shadow partition function converges ($R = 2\pi$,
Borel entire); the topological string diverges
(Gevrey-$1$).  Is $Z_{\mathrm{top}}
= Z^{\mathrm{sh}} \cdot (1 + Z_{\mathrm{np}})$, where
$Z_{\mathrm{np}}$ encodes worldsheet instantons?
\end{openproblem}

\begin{openproblem}[Constructive CY gluing]
\label{op:programme-a-cy-gluing}
\label{op:global-cy-category}
\label{op:programme-h-descent}
\label{conj:ade-chart-k3}
\label{sec:cy-local-global}
\label{rem:db-k3xe-reconstruction}%
\index{Calabi--Yau category!constructive gluing}
$D^b(K3 \times E) \simeq D^b(K3) \otimes D^b(E)$ is
indecomposable (no SOD: Serre functor $[3]$ on a CY$_3$).
The \v{C}ech descent spectral sequence recovers $D^b(K3)$
from two affine charts with $A_\infty$-bimodule transition
data.  At the Kummer point: BKR equivalence
$D^b(\operatorname{Hilb}^2(T^4/\bZ_2)) \simeq
D^b_{\bZ_2}(T^4)$.  Minimum chart number, Brauer
obstruction ($h^{0,2} = 1$), and extension beyond ADE
singularities are open.
\end{openproblem}

\begin{remark}[Master data table]
\label{rem:master-data-k3xe}
\index{K3 x E!master data}
\begin{center}
\small
\renewcommand{\arraystretch}{1.3}
\begin{tabular}{ll}
\textbf{Invariant} & \textbf{Value} \\ \hline
$\dim_\bC$ & $3$ \\
$\chi$ & $0$ \\
$h^{1,1} = h^{2,1}$ & $21$ \\
$\mathrm{rank}\,K_0$ & $48$ \\
$\dim \mathrm{HH}^*$ & $96$ \\
$\kappa(\cA)$ & $3$ \\
$\kappa(K3)$ & $2$ \\
$\kappa(A_E)$ & $24$ \\
$\kappa_{\mathrm{BKM}}$ & $5$ \\
$\kappa_{\mathrm{BCOV}} = \chi/24$ & $0$ \\
Shadow depth & class M ($r_{\max} = \infty$) \\
$F_1 = \kappa/24$ & $1/8$ \\
$F_2$ & $7/1920$ \\
$n^0_1$ & $24 = \chi(K3)$ \\
$S_{\mathrm{BH}}(\Delta)$ & $4\pi\sqrt{\Delta}$ \\
$Z_{\mathrm{DT},0}$ & $1$ \\
$c(A_E)$ & $24$
\end{tabular}
\end{center}
\end{remark}
